% 第7批：《大学物理学学习指导》第5版，第7章“气体动理论基础”，例7.1—7.6。
% 源书印刷页103—105，用户扫描PDF第111—113页。原题条件复核；解答为研究性重写。

\sourceexample{7.1}{103}{111}
\problemtext 氧气和氢气温度相同，已知每个氢气分子的平均平动动能为
$6.21\times10^{-21}\,\mathrm J$。求氧气分子的平均平动动能、氧气方均根速率以及共同的温度。
\workings
对经典理想气体，平均平动动能与分子质量无关，
\[
\bar\varepsilon_{\mathrm{tr}}=\frac32 k_{\mathrm B}T,
\qquad
\boxed{\bar\varepsilon_{\mathrm{tr,O_2}}
=6.21\times10^{-21}\,\mathrm J}.
\]
所以
\[
\boxed{T=\frac{2(6.21\times10^{-21})}{3k_{\mathrm B}}
\approx300\,\mathrm K}.
\]
氧气摩尔质量取$M_{\mathrm{O_2}}=0.032\,\mathrm{kg/mol}$，
\[
\boxed{v_{\mathrm{rms,O_2}}
=\sqrt{\frac{3RT}{M_{\mathrm{O_2}}}}
\approx483\,\mathrm{m/s}}.
\]
\crosscheck 同温不同气体平均平动动能相同，但方均根速率因摩尔质量不同而不同；
$3RT/M_{\rm mol}$单位为$\mathrm{m^2/s^2}$。

\sourceexample{7.2}{103}{111}
\problemtext 容积$V=2\,\mathrm L$的容器装有刚性双原子理想气体，压强$p=1.5\times10^5\,\mathrm{Pa}$。
计算全体分子的平动动能、转动动能、总动能以及气体内能。
\workings
双原子刚性分子有三个平动、两个转动自由度。由$pV=Nk_{\mathrm B}T$，
\[
pV=1.5\times10^5(2\times10^{-3})=300\,\mathrm J.
\]
所以
\[
\boxed{E_{\mathrm{tr}}=\frac32pV=450\,\mathrm J},\qquad
\boxed{E_{\mathrm{rot}}=\frac22pV=300\,\mathrm J},
\]
\[
\boxed{E_{\mathrm{kin}}=\frac52pV=750\,\mathrm J},
\qquad
\boxed{U=750\,\mathrm J}.
\]
\crosscheck 这个模型不计分子相互作用势能与振动自由度；
平均总动能与整体内能相等。

\sourceexample{7.3}{104}{112}
\problemtext 平衡态气体的麦克斯韦速率分布为$f(v)$，分子总数为$N$、最概然速率为$v_p$。
说明$f(v)\,dv$、$\int_{v_p}^\infty Nf(v)\,dv$、
$\int_{v_p}^\infty f(v)\,dv$以及
$\int_0^\infty\tfrac12mv^2f(v)\,dv$各自含义。
\workings
当$f$以$\int_0^\infty f(v)\,dv=1$归一化时，
微小速率区间$[v,v+dv]$内的分子数比例为
\[
\boxed{dP=f(v)\,dv}.
\]
速率不小于最概然速率$v_p$的\textbf{分子数}为
\[
\boxed{N_{v\ge v_p}=\int_{v_p}^\infty Nf(v)\,dv}.
\]
而相同区间内分子数占全部分子的\textbf{比例}为
\[
\boxed{P(v\ge v_p)=\int_{v_p}^\infty f(v)\,dv}.
\]
动能加权平均
\[
\boxed{\int_0^\infty\frac12mv^2 f(v)\,dv
=\left\langle\frac12mv^2\right\rangle
=\frac32k_{\mathrm B}T}.
\]
\crosscheck $f(v)$的单位为$(\mathrm{m/s})^{-1}$，
因此$f(v)dv$无量纲；第四个积分的单位为焦耳。
注意最概然速率不必平分分布曲线下的面积。

\sourceexample{7.4}{104--105}{112--113}
\problemtext 氦气分子的麦克斯韦速率分布曲线峰位为$1000\,\mathrm{m/s}$。
在相同温度下画出（概念上）氢气分布的位置与宽窄趋势，
并求氢气的最概然速率与方均根速率。
\workings
在同温经典理想气体中
\[
v_p=\sqrt{\frac{2RT}{M_{\mathrm{mol}}}},\qquad
v_{\mathrm{rms}}=\sqrt{\frac{3RT}{M_{\mathrm{mol}}}}.
\]
$M_{\mathrm{He}}\simeq4\,\mathrm{g/mol}$，
$M_{\mathrm{H_2}}\simeq2\,\mathrm{g/mol}$，所以
\[
\boxed{v_{p,\mathrm{H_2}}
=\sqrt{\frac{4}{2}}\,v_{p,\mathrm{He}}
=\sqrt2(1000)\approx1.41\times10^3\,\mathrm{m/s}}.
\]
进一步
\[
\boxed{v_{\mathrm{rms,H_2}}
=\sqrt{\frac32}\,v_{p,\mathrm{H_2}}
=\sqrt3(1000)\approx1.73\times10^3\,\mathrm{m/s}}.
\]
概念图上氢气分布在相同温度下向较大速率移动，
曲线更宽、峰高较低，曲线下面积仍为1。
\crosscheck 原书给出$\sqrt2$和$\sqrt3$两项速率倍数，
分别对应氢气最概然速率与方均根速率，不能混为一谈。

\sourceexample{7.5}{105}{113}
\problemtext 显像管中残余空气压强为$10^{-5}\,\mathrm{mmHg}$，温度$27^\circ\mathrm C$。
设空气分子有效直径$d=3.0\times10^{-10}\,\mathrm m$，
摩尔质量$M=28.9\times10^{-3}\,\mathrm{kg/mol}$；
已知$760\,\mathrm{mmHg}=1.013\times10^5\,\mathrm{Pa}$。
求数密度、平均自由程与每分子平均碰撞频率。
\workings
采用题设压强换算和$T\simeq300\,\mathrm K$，
\[
p=\frac{10^{-5}}{760}(1.013\times10^5)\,\mathrm{Pa},
\qquad
\boxed{n=\frac p{k_{\mathrm B}T}\approx3.22\times10^{17}\,\mathrm{m^{-3}}}.
\]
硬球模型给
\[
\boxed{\bar\lambda
=\frac1{\sqrt2\pi d^2 n}
=\frac{k_{\mathrm B}T}{\sqrt2\pi d^2p}
\approx7.8\,\mathrm m}.
\]
平均速率为
\[
\bar v=\sqrt{\frac{8RT}{\pi M}}\approx469\,\mathrm{m/s},
\]
故
\[
\boxed{\bar Z=\frac{\bar v}{\bar\lambda}
\approx60\,\mathrm{s^{-1}}}.
\]
\crosscheck 分子间平均自由程可能超过显像管尺寸；
因此真实粒子与管壁碰撞频繁，但这里按题设求\emph{分子间}平均碰撞频率。

\sourceexample{7.6}{105}{113}
\problemtext 利用玻耳兹曼分布推导等温、静止、均匀重力场中
大气压强随高度$y$的变化规律。
\workings
取分子质量$m$，每分子的重力势能$E_p=mgy$。
玻耳兹曼分布给出数密度
\[
n(y)=n_0\exp\left(-\frac{mgy}{k_{\mathrm B}T}\right).
\]
温度假定处处恒定，局域理想气体状态方程$p=nk_{\mathrm B}T$，
故
\[
\boxed{p(y)=p_0\exp\left(-\frac{mgy}{k_{\mathrm B}T}\right)}
=\boxed{p_0\exp\left(-\frac{M_{\mathrm{mol}}gy}{RT}\right)}.
\]
\crosscheck 指数无量纲；高度增加则压强降低。
真实大气温度随高度变化，此处的单一指数标高律适用于题设等温平衡近似。
